\documentclass[11pt]{article}

% Change "review" to "final" to generate the final (camera-ready) version.
\usepackage[review]{acl}

\usepackage{times}
\usepackage{latexsym}
\usepackage[T1]{fontenc}
\usepackage[utf8]{inputenc}
\usepackage{microtype}
\usepackage{inconsolata}
\usepackage{graphicx}
\usepackage{booktabs}
\usepackage{multirow}
\usepackage{amsmath}
\usepackage{xcolor}

% Placeholder command — bright orange so it is impossible to miss
\newcommand{\placeholder}[1]{\textcolor{orange}{\textbf{[PLACEHOLDER: #1]}}}

\title{Decomposing Automated Fact-Checking:\\
Retrieval, Reranking, and NLI-Guided Classification\\
for Climate Science Claims}

\author{
  \textbf{[Author 1]} \and \textbf{[Author 2]} \and \textbf{[Author 3]} \\
  School of Computing and Information Systems \\
  The University of Melbourne \\
  \texttt{[Group Number]}
}

\begin{document}
\maketitle

% ================================================================
\begin{abstract}
% ================================================================

Automated fact-checking over large evidence corpora requires solving two
interdependent sub-problems: retrieving relevant evidence passages and
classifying a claim given that evidence.
We present a system for verifying climate science claims against a corpus
of 1.2 million evidence passages, and argue that treating these two
components as a unified pipeline obscures where performance is actually
lost.
We propose a diagnostic framework that evaluates retrieval and
classification independently, and use it to guide a four-stage
development process.
Starting from a dense retrieval baseline (Evidence F-score = 0.147,
Accuracy = 52.6\%), we show that hybrid retrieval alone does not improve
precision despite improving recall, motivating a cross-encoder reranker
that raises the Evidence F-score to 0.238 (+65\% relative).
Replacing the classifier initialisation with a model pre-trained on
Natural Language Inference (NLI) further improves classification accuracy
to 59.1\%, particularly on the \textsc{Refutes} (+7.4\%) and
\textsc{Not\_Enough\_Info} (+12.2\%) classes.
Our best system achieves a Harmonic Mean of 0.334 on the development set.
\placeholder{Update abstract with final test-set numbers and System 4 result before submission.}

\end{abstract}

% ================================================================
\section{Introduction}
% ================================================================

The proliferation of unverified claims about climate science has made
automated fact-checking an increasingly important NLP task
\citep{diggelmann-etal-2020-climate}.
Given a natural-language claim, a fact-checking system must (1) retrieve
the most relevant evidence passages from a large corpus, and (2) classify
the claim as \textsc{Supports}, \textsc{Refutes},
\textsc{Not\_Enough\_Info}, or \textsc{Disputed} based on that evidence.

A key challenge that is often underappreciated is that these two
components are tightly coupled: a classifier that receives the wrong
evidence will fail regardless of its capacity, while a retriever that
returns topically similar but logically irrelevant passages creates a
distribution mismatch that degrades classification at inference time.
This coupling makes it difficult to diagnose where a system is failing
from end-to-end metrics alone.

We address this by adopting a \textbf{decomposed evaluation strategy}:
we measure retrieval quality independently (Recall@100, F@5) and
classification quality both with ground-truth evidence (an upper bound)
and with retrieved evidence (the realistic setting).
The gap between these two classification numbers directly quantifies how
much retrieval quality is limiting the classifier.

Our contributions are as follows:
\begin{itemize}
  \item We demonstrate that bi-encoder hybrid retrieval improves recall
        but not precision, and identify the ranking step as the primary
        bottleneck.
  \item We show that a cross-encoder reranker trained on hard negatives
        from the coarse retrieval pool yields the largest single
        improvement in our pipeline (+65\% relative Evidence F-score).
  \item We show that initialising the classifier from an NLI-pretrained
        model substantially improves performance on minority classes,
        particularly \textsc{Refutes} and \textsc{Not\_Enough\_Info}.
  \item \placeholder{Add contribution of System 4 (novel method) here.}
\end{itemize}

% ================================================================
\section{Approach}
% ================================================================

All systems share the same high-level pipeline: a \textbf{retrieval
stage} that selects a small set of evidence passages for each claim,
followed by a \textbf{classification stage} that predicts the claim
label given those passages.
We describe four systems in order of increasing sophistication; each is
motivated by a specific weakness identified in the previous one.

% ----------------------------------------------------------------
\subsection{System 1 — Dense Retrieval Baseline}
\label{sec:system1}
% ----------------------------------------------------------------

\paragraph{Retrieval.}
We encode all 1.2 million evidence passages using
\texttt{all-MiniLM-L6-v2} \citep{reimers-gurevych-2019-sentence}, a
22M-parameter sentence encoder producing 384-dimensional embeddings.
Embeddings are L2-normalised and stored in a matrix; at query time, each
claim is encoded and the top-5 passages are retrieved by cosine
similarity.
Encoding the full corpus takes approximately 3.8 minutes on a single GPU.

\paragraph{Classifier.}
We fine-tune \texttt{DeBERTa-v3-base} \citep{he-etal-2021-deberta} as a
4-class sequence classifier.
The input is the concatenation \texttt{[CLS] claim [SEP] evidence$_1$
$\ldots$ evidence$_k$ [SEP]}, truncated to 512 tokens.
We use a \textbf{two-stage training} procedure motivated by the
distribution mismatch between training and inference:
\begin{enumerate}
  \item \textbf{Stage 1} (10 epochs, lr = 2e-5): train on ground-truth
        evidence, giving the model a clean signal about the relationship
        between claims and correct evidence.
  \item \textbf{Stage 2} (5 epochs, lr = 5e-6): fine-tune on retrieved
        evidence, adapting the model to the noisy, imperfect evidence it
        will see at inference time.
\end{enumerate}
All experiments use fp32 precision; we found that DeBERTa-v3's
disentangled attention mechanism produces NaN gradients under bf16, a
known numerical stability issue with this architecture.

\paragraph{Key finding.}
Despite using a dedicated sentence encoder, dense retrieval achieves
F@5 = 0.147, identical to a TF-IDF baseline.
This suggests that a general-purpose sentence encoder lacks the
domain-specific vocabulary needed to distinguish relevant climate science
evidence from topically similar but logically unrelated passages.

% ----------------------------------------------------------------
\subsection{System 2 — Hybrid Retrieval}
\label{sec:system2}
% ----------------------------------------------------------------

\paragraph{Motivation.}
System 1 showed that dense retrieval alone does not outperform sparse
retrieval.
We hypothesise that the two retrievers are complementary: TF-IDF excels
at exact keyword matches (e.g.\ specific temperature values, chemical
formulae), while dense retrieval captures paraphrases and semantic
equivalences.
Combining them should improve recall without sacrificing precision.

\paragraph{Retrieval.}
We replace \texttt{all-MiniLM-L6-v2} with
\texttt{bge-base-en-v1.5} \citep{xiao-etal-2023-bge}, a 110M-parameter
encoder producing 768-dimensional embeddings that is specifically
optimised for retrieval tasks.
We combine TF-IDF and BGE scores using \textbf{Reciprocal Rank Fusion}
(RRF) \citep{cormack-etal-2009-reciprocal}:
\begin{equation}
  \text{RRF}(d) = \sum_{r \in \{\text{TF-IDF},\, \text{BGE}\}}
  \frac{1}{k + \text{rank}_r(d)}, \quad k = 60
\end{equation}
where $\text{rank}_r(d)$ is the rank of document $d$ under retriever $r$.
We sweep $k \in \{3, 5, 7, 10\}$ on the development set and select the
value that maximises F@5.

\paragraph{Classifier.}
Identical to System 1.

\paragraph{Key finding.}
Hybrid retrieval achieves Recall@100 = 60.7\%, confirming that the
correct evidence is present in the coarse candidate pool for the majority
of claims.
However, F@5 remains flat at 0.144 (Table~\ref{tab:retrieval}).
Analysis reveals why: TF-IDF and BGE each uniquely cover approximately
6\% of ground-truth evidence passages, and their union covers 25\% of
ground truth.
The bottleneck is therefore not recall but \textbf{precision}: both
retrievers score passages by surface-level or semantic similarity to the
claim, but neither can model the logical relationship between a claim and
a passage (e.g.\ whether the passage entails or contradicts the claim).
This motivates a cross-encoder reranker in System 3.

% ----------------------------------------------------------------
% ----------------------------------------------------------------
\subsection{System 3 --- \placeholder{TO BE WRITTEN BY TEAMMATE}}
\label{sec:system3}
% ----------------------------------------------------------------

\placeholder{%
  SYSTEM 3 (final submission system) --- to be written by the teammate
  responsible for the reranker + NLI classifier work.

  Key points to cover:
  \begin{itemize}
    \item Motivation: S2 showed Recall@100 = 60.7\% but F@5 flat at 0.144.
          The coarse pool has the right evidence; the problem is ranking it.
          Bi-encoders cannot model claim--evidence logical interaction.
    \item Retrieval: same hybrid coarse top-100, then DeBERTa cross-encoder
          reranker trained on hard negatives (neg\_ratio=3, 3 epochs,
          max\_length=256). Cite \citet{nogueira-cho-2019-passage}.
    \item Classifier: NLI-pretrained DeBERTa (cross-encoder/nli-deberta-v3-base),
          3-class head replaced with 4-class, same two-stage training.
          Motivation: NLI labels (entailment/contradiction/neutral) map onto
          SUPPORTS/REFUTES/NEI. Cite \citet{nie-etal-2020-adversarial}.
    \item Key finding: F@5 0.144 $\to$ 0.238 (+65\%), accuracy 56.5\% $\to$ 59.1\%.
  \end{itemize}
}

% ----------------------------------------------------------------
\subsection{System 4 --- \placeholder{Novel Method Name}}
\label{sec:system4}
% ----------------------------------------------------------------

\placeholder{%
  SYSTEM 4 (novel method) --- to be written by the teammate responsible
  for the novel contribution.

  Suggested structure:
  \begin{itemize}
    \item Motivation: what specific weakness of System 3 does this address?
    \item Method description with equations/architecture details.
    \item Key design decisions and justification.
    \item Citations to the 1--2 papers that inspired this.
  \end{itemize}
  Target length: roughly 0.5 columns.
}

% ================================================================
\section{Experiments}
% ================================================================

\placeholder{%
  EXPERIMENTS section --- to be written collaboratively.

  \subsection*{Experimental Setup}
  Cover: dataset stats (Table of label distribution already drafted above),
  evaluation metrics (F, A, H), the diagnostic metric A(GT) vs A(Retrieved),
  hardware, random seeds, training times.

  \subsection*{Retrieval Results}
  Table comparing Recall@100 and F@5 across all four systems.
  Template already drafted --- fill in System 4 numbers.

  \subsection*{End-to-End Results}
  Main results table across all four systems.
  Include A(GT) column --- this is the key diagnostic number.
  Template already drafted --- fill in missing cells.
}

% ================================================================
\section{Analysis}
% ================================================================

\placeholder{%
  ANALYSIS section --- to be written collaboratively.

  \subsection*{Retrieval is the Primary Bottleneck}
  Use the A(GT) vs A(Retrieved) gap to quantify how much retrieval
  quality limits the classifier. This is the central argument.

  \subsection*{Why Bi-Encoders Fail at Precision}
  Explain the 65\% relative F@5 improvement from the reranker.
  Bi-encoders score independently; cross-encoders read jointly.

  \subsection*{NLI Pre-training and Minority Classes}
  REFUTES +7.4\%, NEI +12.2\% from NLI initialisation.
  Explain the structural alignment between NLI and fact-checking labels.
  Note DISPUTED remains 0\% --- discuss why.

  \subsection*{Qualitative Analysis}
  % --- QUALITATIVE PLACEHOLDER ---
  % 2--3 case studies drawn from the dev set predictions.
  % To find candidates: cross-reference dev-claims-predictions-v2.json,
  % dev-claims-predictions-v3.json, and dev-claims-predictions-v7.json
  % against dev-claims.json.
  %
  % CASE STUDY 1 --- Retrieval failure (motivates the reranker):
  %   Find a claim where V2/V3 retrieved wrong evidence but V4/V7 got it right.
  %   Show: claim text | GT evidence (key sentence) | retrieved evidence (key sentence)
  %   Show: S1 prediction (wrong) vs S3 prediction (correct)
  %   Argument: the classifier is not at fault; the retriever is.
  %
  % CASE STUDY 2 --- Classification failure despite correct retrieval:
  %   Find a claim where all systems retrieved correct evidence but label is still wrong.
  %   Show: claim text | evidence | why the logical relationship is subtle.
  %   Argument: retrieval is not the only bottleneck.
  %
  % CASE STUDY 3 (optional) --- DISPUTED failure:
  %   Pick any DISPUTED claim. Show evidence exists on both sides.
  %   Argument: task requires detecting evidential conflict, not just relevance.
  %
  % Format each as a small indented block:
  %   Claim: ``...''
  %   GT evidence: ``...'' (evidence-XXXXX)
  %   Retrieved (S1): ``...'' (evidence-YYYYY)
  %   S1 pred: X  |  S3 pred: Y  |  Ground truth: Z
  % Then 2--3 sentences of analysis. Target: ~0.5 page total.
  % --- END QUALITATIVE PLACEHOLDER ---
}

% ================================================================
\section{Conclusion}
% ================================================================

\placeholder{%
  CONCLUSION --- to be written after all results are in.

  Cover:
  - Retrieval quality is the dominant bottleneck (quantify with final numbers).
  - Cross-encoder reranking is the single most impactful improvement.
  - NLI pre-training transfers well to fact verification.
  - DISPUTED class remains unsolved --- structural challenge.
  - Limitations: Colab compute, small training set, class imbalance.
  - Future work: iterative retrieval, larger reranker, better DISPUTED handling.
  - Add 1--2 sentences on System 4's contribution.
  Target: ~0.25 page.
}

% ================================================================
\section*{Limitations}
% ================================================================

\placeholder{%
  LIMITATIONS --- brief, ~1 paragraph.
  Mention: fp32-only constraint, 1228 training examples, DISPUTED class,
  Colab compute ceiling, test-set distribution shift.
}

% ================================================================
\section*{Team Contributions}
% ================================================================

\placeholder{%
  Fill in before submission. Example format:
  [Author 1]: Systems 1 and 2 (pipeline\_v2.py, pipeline\_v3.py),
              Abstract, Approach sections 2.1--2.2, Experiments setup.
  [Author 2]: Systems 3 and 4, Approach sections 2.3--2.4, Analysis.
  [Author 3]: Report coordination, Results tables, Conclusion,
              qualitative analysis, leaderboard submissions.
}

% ================================================================
\bibliography{custom}
% ================================================================

\end{document}
